\documentclass[aspectratio=169]{beamer}

% Theme and basic setup
\usetheme{Madrid}
\usecolortheme{default}
\setbeamertemplate{navigation symbols}{}

\usepackage[utf8]{inputenc}
\usepackage[T1]{fontenc}
\usepackage{inconsolata} % nicer monospace
\usepackage{amsmath,amssymb}
\usepackage{mathtools}
\usepackage{listings}
\usepackage{xcolor}

\input{unicodeletters}

% Listings setup (approximate Agda using Haskell-like settings)
\lstdefinelanguage{Agda}{
  morekeywords={data,record,module,where,open,import,let,in,
    Set,Set₁,Set₂,postulate,with,pattern,mutual,infix,infixl,infixr,
    forall,constructor,field,boxed,unquote,quote,unquoteDecl},
  sensitive=true,
  morecomment=[l]--,
  morecomment=[s]{\{-}{-\}},
  morestring=[b]"
}

\lstset{
  language=Agda,
  basicstyle=\ttfamily\small,
  keywordstyle=\bfseries\color{blue!70!black},
  commentstyle=\itshape\color{gray!70!black},
  stringstyle=\color{green!50!black},
  showstringspaces=false,
  columns=fullflexible,
  keepspaces=true,
  frame=single,
  breaklines=true,
  tabsize=2,
  literate=
  {≡}{{$\equiv$}}1
  {->}{{$\to$}}1
}

% Title info
\title[Agda by Example]{Agda by Example:\\Programming and Proving with Dependent Types}
\author{Peter Thiemann}
\institute{Freiburg University}
\date{March 13, 2026}

\begin{document}

%------------------------------------------------------------
\begin{frame}
  \titlepage
\end{frame}

%------------------------------------------------------------
\begin{frame}{Motivation}
  \begin{itemize}
    \item We want software that is \textbf{correct by construction}, not “probably ok after tests”.
    \item Types as \textbf{specifications}: function signatures as contracts.
    \item Dependent types let types talk about \textbf{values}: lengths, bounds, invariants.
    \item Agda fulfills two roles:
      \begin{itemize}
        \item a \textbf{total functional programming language}, and
        \item an \textbf{interactive proof assistant}.
      \end{itemize}
  \end{itemize}
\end{frame}

%------------------------------------------------------------
\begin{frame}{Agda in a Nutshell}
\small
\begin{itemize}
  \item Agda is a proof assistant and dependently typed FP language based on
  \textbf{intuitionistic Martin--L\"of type theory} (constructive proofs-as-programs).

  \item Agda is named after the Swedish song "Hönan Agda" about a
    hen, serving as a pun on the proof assistant Coq (another
    dependently typed proof assistant of French origin).

  \item Agda was primarily developed by \textbf{Ulf Norell}
    (notably in his 2007 PhD work), taking inspiration from ALF,
    then grown by a broader community into today’s widely used
    system. 

  \item Implemented in \textbf{Haskell}; emphasizes
    \textbf{interactive development} (holes), and \textbf{dependent pattern matching}.

  \item The \textbf{Agda standard library} provides reusable foundations
  (data structures, algebra, equality reasoning) and strongly shapes idiomatic Agda.

  \item Good for \textbf{mechanized metatheory}, \textbf{certified programming},
    and teaching the tight link between \textbf{types, specs, and proofs}.
\end{itemize}
\end{frame}

%------------------------------------------------------------
\begin{frame}{What We Will Do (90 min)}
  \begin{enumerate}
    \item Warm-up: simple programs, total functions.
    \item Inductive data types: lists, recursion.
    \item Dependent types in action: vectors indexed by length.
    \item Propositions as types: equality and proofs.
    \item Case study: a small verified function.
  \end{enumerate}
  \vfill
  \textbf{Goal:} Leave with a mental model for using types as specifications,
  and how to explore that space using Agda.
\end{frame}

%------------------------------------------------------------
\section{Getting Started in Agda}

\begin{frame}{Agda as a Language + Proof Assistant}
  \begin{itemize}
    \item Source files: \texttt{.agda} or literate \texttt{.lagda.md}.
    \item Interactive editing:
      \begin{itemize}
        \item Emacs, VSCode, etc.
        \item Holes: typed as \texttt{?}, typeset as
          \texttt{\{\!~\!\}}; most commands fill/refine holes.
      \end{itemize}
    \item Everything is total:
      \begin{itemize}
        \item Functions are defined by exhaustive pattern matching.
        \item Termination is checked.
      \end{itemize}
  \end{itemize}
\end{frame}

%------------------------------------------------------------
\begin{frame}[fragile]{Warm-Up: Simple Functions}
  \begin{columns}
    \begin{column}{0.5\textwidth}
      \lstinline!module Warmup where!

\begin{lstlisting}
open import Agda.Builtin.Nat

double : Nat -> Nat
double n = n + n

triple : Nat -> Nat
triple n = {!!}
\end{lstlisting}
    \end{column}
    \begin{column}{0.5\textwidth}
      \textbf{Exercise 1} (few minutes):
      \begin{itemize}
        \item Fill the hole for \texttt{triple}.
        \item Check the types with \texttt{C-c C-,}.
        \item Refine the hole with \texttt{C-c C-r}.
        \item Experiment with \texttt{C-c C-n} to normalize expressions.
      \end{itemize}
      \vspace{1em}
      \textbf{Key idea:} Types as a starting point; holes as “typed to-do items”.
    \end{column}
  \end{columns}
\end{frame}

%------------------------------------------------------------
\section{Inductive Data Types}

\begin{frame}[fragile]{Lists from Scratch}
  \begin{columns}
    \begin{column}{0.5\textwidth}
\begin{lstlisting}
data List (A : Set) : Set where
  []   : List A
  _::_ : A -> List A -> List A

length : {A : Set} -> List A -> Nat
length []       = 0
length (x :: xs) = 1 + length xs
\end{lstlisting}
    \end{column}
    \begin{column}{0.5\textwidth}
      \textbf{Key points:}
      \begin{itemize}
        \item Inductive data = constructors + recursion.
        \item Types track element type (\texttt{A}), not size (yet).
        \item Pattern matching drives computation.
      \end{itemize}
    \end{column}
  \end{columns}
\end{frame}

%------------------------------------------------------------
\begin{frame}[fragile]{Appending Lists}
\begin{lstlisting}
append : {A : Set} -> List A -> List A -> List A
append []       ys = ys
append (x :: xs) ys = x :: append xs ys
\end{lstlisting}

  \textbf{Exercise 2:}
  \begin{itemize}
    \item Implement \texttt{reverse : \{A : Set\} -> List A -> List A}.
    \item Use pattern matching and recursion; load the file and
          let Agda check totality and termination.
  \end{itemize}

  \vfill
  \textbf{Observation:}
  \begin{itemize}
    \item Types specify \emph{shape} of data, but not \emph{properties} like length.
  \end{itemize}
\end{frame}

%------------------------------------------------------------
\section{Dependent Types in Action}

\begin{frame}[fragile]{Vectors: Lists with Length in the Type}
  \begin{columns}
    \begin{column}{0.5\textwidth}
\begin{lstlisting}
open import Agda.Builtin.Nat

data Vec (A : Set) : Nat -> Set where
  []   : Vec A 0
  _::_ : {n : Nat} ->
         A -> Vec A n -> Vec A (suc n)
\end{lstlisting}
    \end{column}
    \begin{column}{0.5\textwidth}
      \textbf{Idea:}
      \begin{itemize}
        \item \texttt{Vec A n} is a list of \texttt{A} \emph{indexed} by length \texttt{n}.
        \item Indices (here: \texttt{Nat}) are values appearing in types.
        \item Invariants like “length is preserved” become type-level facts.
      \end{itemize}
    \end{column}
  \end{columns}
\end{frame}

%------------------------------------------------------------
\begin{frame}[fragile]{Mapping over Vectors}
\begin{lstlisting}
vmap :
  {A B : Set} {n : Nat} ->
  (A -> B) -> Vec A n -> Vec B n
vmap f []        = []
vmap f (x :: xs) = f x :: vmap f xs
\end{lstlisting}

  \textbf{Reading the type:}
  \begin{itemize}
    \item For any length \texttt{n}, if you give me a vector of \texttt{A} of length \texttt{n},
    \item I return a vector of \texttt{B} of the \emph{same} length \texttt{n}.
  \end{itemize}

  \vfill
  \textbf{Exercise 3:}
  \begin{itemize}
    \item Define \texttt{vappend : Vec A m -> Vec A n -> Vec A (m + n)}.
  \end{itemize}
\end{frame}

%------------------------------------------------------------
\section{Propositions as Types}

\begin{frame}[fragile]{Equality as a Type}
  In Agda, propositional equality is a data type:

\begin{lstlisting}
data _≡_ {A : Set} (x : A) : A -> Set where
  refl : x ≡ x
\end{lstlisting}

  \begin{itemize}
    \item A value of type \texttt{x ≡ y} is evidence that \texttt{x} and \texttt{y} are equal.
    \item There is only one constructor: \texttt{refl}, proving \texttt{x ≡ x}.
    \item To prove \texttt{x ≡ y}, we typically normalize both sides and use \texttt{refl}.
  \end{itemize}
\end{frame}

%------------------------------------------------------------
\begin{frame}[fragile]{A Simple Proof: Right Identity of +}
\begin{lstlisting}
open import Agda.Builtin.Nat
open import Agda.Builtin.Equality

zero-right-id : (n : Nat) -> n + 0 ≡ n
zero-right-id zero    = refl
zero-right-id (suc n) = cong suc (zero-right-id n)
\end{lstlisting}

  \textbf{Notes:}
  \begin{itemize}
    \item \texttt{cong} lifts equality through a function:
\begin{lstlisting}
cong : {A B : Set} (f : A -> B) {x y : A} ->
       x ≡ y -> f x ≡ f y
\end{lstlisting}
    \item The proof is just a recursive function following the
      structure of \texttt{Nat}.
    \item It constitutes a proof by mathematical induction that $\forall
      n. n + 0 = n$.
  \end{itemize}
\end{frame}

%------------------------------------------------------------
\begin{frame}[fragile]{List Property: Appending []}
\begin{columns}
  \begin{column}{0.5\textwidth}
\begin{lstlisting}
length-append-[] :
  {A : Set} (xs : List A) ->
  length (append xs []) ≡ length xs
length-append-[] [] = ?
length-append-[] (x :: xs) = ?
\end{lstlisting}
    
  \end{column}
  \begin{column}{0.5\textwidth}
    \textbf{Exercise 4:}
    \begin{itemize}
      \item Type this lemma and ask Agda for goals (holes).
      \item Fill in the recursive case using \texttt{cong}.
    \end{itemize}
    \vspace{1em}
    \textbf{Takeaway:}
    \begin{itemize}
      \item Proofs look like ordinary recursive functions.
      \item The type encodes the property.
    \end{itemize}
  \end{column}
\end{columns}
\end{frame}

%------------------------------------------------------------
\begin{frame}[fragile]{List Property: Left and Right Id}
\begin{columns}
  \begin{column}{0.5\textwidth}
\begin{lstlisting}
empty-right-id :
  {A : Set} (xs : List A) ->
  append xs [] ≡ xs
empty-right-id = ?

empty-left-id :
  {A : Set} (xs : List A) ->
  append [] xs ≡ xs
empty-left-id = ?
\end{lstlisting}
    
  \end{column}
  \begin{column}{0.5\textwidth}
    \textbf{Exercise 4a:}
    \begin{itemize}
      \item Type the empty-right lemma and ask Agda for goals (holes).
      \item Fill in the recursive case.
    \end{itemize}
    \vspace{1em}
    \textbf{Exercise 4b:}
    \begin{itemize}
      \item Type the empty-left lemma and ask Agda for goals (holes).
      \item Do you need induction?
    \end{itemize}
    \vspace{1em}
    \pause
    \textbf{Takeaway:}
    \begin{itemize}
      \item Some properties hold \emph{definitionally}, i.e., by normalizing according to function definitions.
    \end{itemize}
  \end{column}
\end{columns}
\end{frame}

%------------------------------------------------------------
\section{Case Study: Verified Lookup}

\begin{frame}[fragile]{Safe Indexing with \texttt{Fin}}
  \begin{columns}
    \begin{column}{0.5\textwidth}
\begin{lstlisting}
data Fin : Nat -> Set where
  fzero : {n : Nat} -> Fin (suc n)
  fsuc  : {n : Nat} -> Fin n -> Fin (suc n)
\end{lstlisting}
    \end{column}
    \begin{column}{0.5\textwidth}
      \textbf{Idea:}
      \begin{itemize}
        \item \texttt{Fin n} represents numbers \texttt{0 .. n-1}.
        \item This range corresponds to the domain of  \texttt{Vec A n}!
        \item Combining \texttt{Fin n} with \texttt{Vec A n} yields safe lookup.
        \item It is impossible to construct an out-of-bounds index. 
      \end{itemize}
    \end{column}
  \end{columns}
\end{frame}

%------------------------------------------------------------
\begin{frame}[fragile]{Lookup That Cannot Fail}
\begin{lstlisting}
lookup :
  {A : Set} {n : Nat} ->
  Fin n -> Vec A n -> A
lookup fzero    (x :: xs) = x
lookup (fsuc i) (x :: xs) = lookup i xs
\end{lstlisting}

  \textbf{Key observations:}
  \begin{itemize}
    \item No case for \texttt{[]} is needed: \texttt{Fin 0} is empty
      (recognized by Agda).
    \item No runtime bounds checking.
    \item The type guarantees safety by construction.
  \end{itemize}

  \vfill
  \textbf{Exercise 5:}
  \begin{itemize}
    \item Implement a small function using \texttt{lookup}
          (e.g., get the head of a non-empty vector).
  \end{itemize}
\end{frame}

%------------------------------------------------------------
\begin{frame}[fragile,fragile]
  \frametitle{Challenge Problem}
\begin{lstlisting}
vsnoc : {A : Set} {n : Nat} -> Vec A n -> A -> Vec A (suc n)
vsnoc xs x = {!!}

vreverse : {A : Set} {n : Nat} -> Vec A n -> Vec A n
vreverse [] = []
vreverse (x :: xs) = vsnoc (vreverse xs) x
\end{lstlisting}
  \vfill
  \textbf{Exercise 6 (challenge):}
  \begin{itemize}
    \item Try to implement a function \texttt{vsnoc} to attach a single
      element to the \emph{end} of a vector. This function can be use
      to implement \texttt{vreverse} to  reverse a vector.
    \item Analyze the problem you run into
  \end{itemize}
\end{frame}
%------------------------------------------------------------
\section{Wrap-Up}

\begin{frame}{What We Saw}
  \begin{itemize}
    \item \textbf{Inductive data} for structure (Nat, List, Vec).
    \item \textbf{Dependent types} to encode invariants (length, bounds).
    \item \textbf{Equality as a type} and small composable proofs.
    \item \textbf{Safe APIs} by design:
      \begin{itemize}
        \item \texttt{Vec A n}, \texttt{Fin n}, \texttt{lookup}.
      \end{itemize}
  \end{itemize}
\end{frame}

%------------------------------------------------------------
\begin{frame}{Takeaways for Day-to-Day Work}
  After this tutorial, you should be able to:
  \begin{itemize}
    \item Encode key invariants in types to design safer APIs.
    \item Formulate precise specifications as type signatures.
    \item Apply proof-driven reasoning to structure and validate behavior.
  \end{itemize}
\end{frame}

%------------------------------------------------------------
\begin{frame}{Where to Go Next}
  \begin{itemize}
    \item Agda standard library documentation.
    \item “Programming Language Foundations in Agda” (PLFA).
    \item Explore:
      \begin{itemize}
        \item More advanced dependently typed programming patterns.
        \item Larger verification case studies: interpreters, DSLs, protocols.
      \end{itemize}
  \end{itemize}
  \vfill
  \centering
  \textit{Thank you!}
\end{frame}

%------------------------------------------------------------
\end{document}
